\section{Simulating Gaussian Random Variables}

To simulate \(Z_{1:n} = (Z_1,\dots,Z_n)^T \sim \normal(0, \Sigma_n)\) one typically
searches for \(\Lmatrix_n\) such that
\[
	\Lmatrix_n \Lmatrix_n^T = \Sigma_n,
\]
because for \(Y_{1:n}\sim\normal(0,\identity)\) independent standard normal
random variables we have
\[
	Z_{1:n} = \Lmatrix_n Y_{1:n} \sim \normal(0, \Sigma_n).
\]
The algorithm of choice is typically the cholesky decomposition.
In particular for a centered gaussian process/random field 
\(Z=(Z(x))_{x\in\real^\dimension}\) with covariance function \(\C\),
we can simulate \((Z(x_1),\dots, Z(x_n))\) by calculating \(\Sigma_n\) from its
covariance function
\[
	\Sigma_n = \begin{pmatrix}
		\C(x_1, x_1) &\dots &\C(x_1, x_n)\\
		\vdots &  & \vdots \\
		\C(x_n, x_1) &\dots &\C(x_n, x_n)
	\end{pmatrix}.
\]
But what if you already have simulated \((Z(x_1),\dots,Z(x_n))\) and now want
to obtain \(Z_{n+1} = Z(x_{n+1})\) without changing your previous values? If you wanted
to do optimization on a random field for example, the position \(X_{n+1}\) might
even be a random variable depending on the previous results. You can hardly
start over in this case.
